\section{Background}

\subsection{High Level Synthesis}
HLS offers a promising way to design hardware at a high-level abstraction such as C-like language, which can release users from designing at register transfer level. HLS compilers~\cite{VivadoHls, legup, Dynamatic_Josipovi, uIR_Sharifian} then automatically convert the high-level language to HDLs by three major processes: allocation, scheduling, and binding. Allocation places all compute units on the datapath, and binding determines where each operation executes.

% \hl{first gives an overview of high level synthesis.}
\begin{figure}[t]
  \centering
  \includegraphics[width=\linewidth]{figure/background.pdf}
  \caption{Static and dynamic scheduling in high level synthesis.}
  \label{fig:background}
  \vspace{-9pt}
\end{figure}

Scheduling is the most important step in modern HLS tools. There are two different scheduling methodologies: static and dynamic, as shown in \autoref{fig:background}. The most common flow of static approach consists of 4 steps~\cite{intro-hls}: software compilation, static scheduling, allocating \& binding, and finite state machine (FSM) construction, as shown in \autoref{fig:background}(a). Compilation transforms the high-level program into a software intermediate representation. Then, static scheduling algorithm determines the execution time for operations considering dependencies and resource constraints. For example, the add and shift-right operations are placed in the first cycle (\verb|C0|). Allocating \& binding algorithms explore the opportunity for resource sharing, allocate necessary resources, and bind operations to them. Finally, a finite state machine is constructed for RTL generation. 

Static scheduling~\cite{rl_hls, cdl_hls, thread_weaving} is very effective for statically predictable programs such as perfect loops.  Pipelining is a key optimization technique to exploit parallelism among multiple loop iterations. In the sequential execution, one iteration can only begin after the previous iteration is finished. Pipelining allows iterations to overlap which improves throughput and gain resource sharing. The distance between two adjacent iterations is called initial interval (II)~\cite{}, which is the indicator of throughput.

% \hl{it seems that dynamic scheduling is not used in the following parts.}
Dynamic scheduling generates a dataflow circuit by leveraging elastic units~\cite{elastic_circuit}, such as merge, branch, etc., as shown in \autoref{fig:background}(b). All data signals in dataflow circuits are accompanied by handshake signals, which are \textit{valid} and \textit{ready}, in opposite directions, indicating the availability of the next data from the source unit and the readiness of the target unit to accept it, respectively. Static scheduling makes conservative assumptions for unresolvable dependencies at compile time. Dynamic scheduling overcomes this inefficiency by postponing the scheduling decisions until run time, However, dynamic scheduling suffers from high resource consumption. Hybrid scheduling combines static and dynamic schedulings~\cite{cheng_dss, find_islands}.

\subsection{Intermediate Representation}
Intermediate representation is an important abstraction to simplify the design of compilation and synthesis flows. Compiler infrastructures like LLVM~\cite{lattner2004llvm} and GCC~\cite{novillo2003tree} apply machine-independent optimizations and generate code for different target architectures based on IRs. For example, Static single assignment (SSA) form~\cite{StaticSingleAssignment} and control dataflow graph (CDFG) representation~\cite{cdfg_allen} are widely used in compiler optimizations and static analysis. Many HLS tools~\cite{VivadoHls, legup, Dynamatic_Josipovi,scalehls,scalehls} adopt LLVM IR as their internal representation. However, LLVM is a pure software IR that doesn't contain any hardware information. \cite{LLHD_Schuiki, firrtl, ahir,uIR_Sharifian, Calyx_Rachit} present hardware IRs that provide low-level abstraction to support hardware design. \hl{let us add liqiang's TENET paper into the list.} 

%Calyx~\cite{Calyx_Rachit} combines software-like control flow and structural representation to describe architectural design.
%Hardware IR are used in hardware synthesis flows.

%Halide~\cite{Halide} aims at explore parallelism and locality of software domains like image processing.
% Relay IR~\cite{relay} proposes a differentiable language for expressing machine learning models. 

MLIR~\cite{mlir} is a novel compiler infrastructure that provides powerful scalability and modularity. MLIR greatly facilitates the implementation of various IRs and transformations among them. All IRs in MLIR follow the SSA form and an explicit type system. \emph{Dialect} in MLIR is a hierarchical structure template allowing basic optimizations in MLIR to be reused. This hierarchy supplies an expressive representation, which makes it easy to implement a flexible IR. 

% \section{Background}
% \subsection{HLS Scheduling}
% Traditional HLS tools like Vivado HLS\cite{VivadoHls} and LegUp\cite{legup} adopt static scheduling like SDC algorithm\cite{cong_sdc} that decides certain clock cycle in the compile time, allocates each operation to a state and uses Finite State Machine (FSM) to deal with control flow. Therefore, hardware implementation works in a sequential way. To increase parallelism, explicit pipeline optimizations like module scheduling~\cite{rau_module_scheduling} are adopted. When pipelining the program, static scheduling needs to make a conservative scheduling at the compile time that reserves enough time to ensure all dependence and memory conflicts cannot be violated. In a word, if the program has complicated control flow and the branches are imbalanced, these HLS tools may choose a big Initial Interval (II) which leads to a bad performance. One benefit of static scheduling is that we can analysis whether every pair of operations can share the same resource, and gain enough resource sharing.

% To overcome the conservative pipeline in the compile time, Dynamatic~\cite{Dynamatic_Josipovi} adopts elastic circuits techniques\cite{elastic_circuit} in dynamic scheduling. Each operation in CDFG is wrapped to a pipeline-style compute unit, and the control flow is converted to the combination of many elastic components such as fork, merge, branch and select. As for memory conflicts, their work adopts a Load Store Queue (LSQ)~\cite{load_store_queue} to dispute load and store operations. Every pair of operations that has dependence in CDFG is communicating through a handshake protocol\cite{handshake}. By the use of these signals, we can satisfy the dependence during the run-time and overcome the shortcoming of conservative selection. Lacking enough information at the compile time, resource sharing between different operations cannot be obtained. To get a trade-off between performance and resource utilization, DSS~\cite{cheng_dss} first combines these two scheduling method in functional level by manually determining whether each function uses dynamic scheduling. The static parts are then wrapped with handshake protocol, which are directly connected to dynamic parts at RTL level.

% \subsection{Intermediate Representation in HLS}
% Intermediate representation is an important abstraction that represents a model of software or hardware, and can be applied on compiler transformation for analysis and optimizations. As a powerful software compiler infrastructure, LLVM is often used in HLS tools like Vivado HLS and LegUp. Because LLVM IR doesn't aim at hardware, some inner representations are extra implemented in those tools. Existing HLS tools don't consider about readability and flexibility, for which these inner representations are not exposed to users.

% \subsection{MLIR Infrastructure}
% MLIR\cite{MLIR} is a novel compiler infrastructure which provides powerful scalability and modularity, and is particularly suitable for multi-level IR transformation. The \emph{Dialect} in MLIR is a set of operations, attributes and types, which can be regarded as a flexible IR. All the IRs in MLIR obey Static Single Assignment (SSA) form~\cite{StaticSingleAssignment} and explicit type system. MLIR provides several built-in Dialects, such as \emph{Affine}, \emph{SCF} and \emph{Standard}. \emph{Affine} provides a abstraction for affine operations like affine.for, which can be used to polyhedral compilation techniques. \emph{SCF} describes static control flow in a higher-level abstraction keeping branch and loop structure, which is helpful for analyzing and optimizing. \emph{Standard} is a collection of basic operation such as addition, comparison and function call. These dialects can be used to describe a program at the same time, which is also the feature provided by MLIR, called progressive lowering. The program described in \emph{Affine}, \emph{SCF} and \emph{Std} has similar functionality and representation with high-level language.

% \subsection{Motivation}
% Most of HLS tools only adopts static scheduling, while several tools such as Dynamatic and SOFF\cite{SOFF} only considers about dynamic scheduling. DSS proposes a mixed scheduling for performance and resource trade-offs. However, the mixed scheduling method that DSS designs is just the simple combination of dynamic and static parts. The method proposed by DSS can be seen as a function-level mixed scheduling approach, which needs an elaborate design to partition an application into multiple kernels. Then, dynamic and static parts are separately sent to Dynamatic and Vivado HLS, which are finally lowered to VHDL. This scheduling method isolates two kinds of scheduling algorithms, which makes the hardware generation separate and independent.

% High-level synthesis substantially reduces the difficulty on hardware design by automatically converting high-level programs to HDLs. But traditional HLS frameworks lack explicit intermediate representation, which makes it hard to inspect and interfere during the whole HLS process. It is the only way to manipulate the hardware generation that users manually insert proper directives like loop unroll and pipeline. Without the information of each step, we can only get the final HDL implementation which is not friendly for understanding and further optimization. In another word, the area and performance of final hardware can only count on HLS compilers, which may generate unpredictable hardware with different behaviors from what we expect. Different HLS tools support different kinds of directives, and may use different scheduling algorithm which cannot be interfered or selected. Explicit IR can provide significant convenience to generate predictable hardware especially during the scheduling and binding phases. For example, users can select different scheduling algorithm provided by our flow, and they can even implement their own scheduling algorithm to generate better hardware with more knowledge of their own applications and reuse the other part in our flow. What's more, some optimizations like resource sharing and memory partition are easily applied, which can also be embodied in this IR.

% So we are willing to implement a HLS framework based on multi-level IR paradigm that is suitable for representing hardware from different abstraction levels and that is also easy for compiler to optimize. To make use of both static and dynamic scheduling, the IR we design should also be capable of representing both dynamic and static scheduling. An IR designed for HLS lessens the gap between high level programs and HDLs, and also provide the ability to debug and manipulate at different level.

% Hardware generated by Vivado HLS can only be interfered by manually inserting directives, while DSS only supports the schedule selection of each function. In contrast, the hardware can be interfered at different level like scheduling and binding strategies with the help of explicit IR.